\documentclass[aspectratio=169]{beamer}

\usepackage[T1]{fontenc}
\usepackage[utf8]{inputenc}
\usepackage[french]{babel}
\usepackage{lmodern}
\usepackage{tikz}
\usepackage{booktabs}
\usepackage{array}
\usepackage{ragged2e}

\usetikzlibrary{arrows.meta,positioning,shapes.geometric}

\definecolor{Primary}{RGB}{19,138,131}
\definecolor{Accent}{RGB}{231,111,81}
\definecolor{Dark}{RGB}{34,40,49}
\definecolor{Soft}{RGB}{245,247,250}
\definecolor{Muted}{RGB}{108,117,125}

\setbeamertemplate{navigation symbols}{}
\setbeamercolor{background canvas}{bg=white}
\setbeamercolor{normal text}{fg=Dark,bg=white}
\setbeamercolor{title}{fg=white}
\setbeamercolor{frametitle}{fg=Dark,bg=Soft}
\setbeamercolor{structure}{fg=Primary}
\setbeamercolor{block title}{fg=white,bg=Primary}
\setbeamercolor{block body}{fg=Dark,bg=Soft}
\setbeamercolor{alerted text}{fg=Accent}
\setbeamerfont{title}{series=\bfseries,size=\Large}
\setbeamerfont{frametitle}{series=\bfseries,size=\large}

\setbeamertemplate{footline}{
  \leavevmode%
  \hbox{%
    \begin{beamercolorbox}[wd=.78\paperwidth,ht=2.8ex,dp=1.2ex,leftskip=1em]{author in head/foot}%
      \color{Muted}\footnotesize micro:bit Cutebot Pro 
    \end{beamercolorbox}%
    \begin{beamercolorbox}[wd=.22\paperwidth,ht=2.8ex,dp=1.2ex,rightskip=1em plus1fil]{date in head/foot}%
      \color{Muted}\footnotesize \insertframenumber{} / \inserttotalframenumber
    \end{beamercolorbox}%
  }%
}

\title{micro:bit Cutebot Pro - Tilt-controlled robot in C}
\author{Paul-henri DJOKO \\ Franck ulrich KENFACK \\ Carlos MENESES}
\institute{ENSTA -- IoT}
\date{2026}

\begin{document}

{
\setbeamertemplate{footline}{}
\begin{frame}[plain]
  \begin{tikzpicture}[remember picture,overlay]
    \fill[Primary] (current page.north west) rectangle ([yshift=-2.8cm]current page.north east);
    \fill[Accent] ([xshift=-2.2cm,yshift=0.6cm]current page.south east) circle (1.7cm);
    \fill[Primary!15] ([xshift=1.2cm,yshift=0.9cm]current page.south west) circle (1.1cm);
  \end{tikzpicture}

  \vspace{0.9cm}
  {\usebeamerfont{title}\color{white}\inserttitle\par}
  \vspace{0.2cm}
  {\large\color{white!90}\insertsubtitle\par}

  \vspace{1.8cm}
  \begin{beamercolorbox}[rounded=true,shadow=false,wd=\textwidth,sep=1.1em]{block body}
    \textbf{Id\'ee centrale :} piloter un robot via une micro:bit de mani\`ere intuitive, interactive et s\'ecuris\'ee, en s'appuyant sur une architecture C lisible.
  \end{beamercolorbox}

  \vfill
  {\small \insertauthor \hfill \insertinstitute}
\end{frame}
}

\begin{frame}{Probl\'ematique}
  \begin{columns}[T,onlytextwidth]
    \column{0.55\textwidth}
    \begin{block}{Question de d\'epart}
      Comment contr\^oler un robot mobile \textbf{sans fil}, de fa\c{c}on
      \textbf{simple}, \textbf{r\'eactive} et \textbf{s\^ure}, avec deux
      micro:bit v2 et sans utiliser MakeCode ?
    \end{block}

    \begin{block}{Objectifs}
      \begin{itemize}
        \item utiliser l'inclinaison comme interface naturelle;
        \item transmettre les commandes par radio en temps r\'eel;
        \item prot\'eger le robot face aux obstacles;
        \item structurer le code pour faciliter les tests et l'\'evolution.
      \end{itemize}
    \end{block}

    \column{0.4\textwidth}
    \begin{alertblock}{Contraintes}
      \begin{itemize}
        \item programmation bas niveau en C;
        \item calibration de l'acc\'el\'erom\`etre;
        \item coordination de plusieurs p\'eriph\'eriques;
        \item m\^eme projet SES pour les deux cartes.
      \end{itemize}
    \end{alertblock}
  \end{columns}
\end{frame}

\begin{frame}{Pr\'esentation du projet}
  \begin{columns}[T,onlytextwidth]
    \column{0.48\textwidth}
    \begin{block}{Partie manette}
      \begin{itemize}
        \item lecture de l'acc\'el\'erom\`etre LSM303AGR;
        \item conversion des valeurs brutes vers l'intervalle `[-100,100]';
        \item lecture des boutons A et B;
        \item envoi des donn\'ees par \texttt{NRF\_RADIO}.
      \end{itemize}
    \end{block}

    \column{0.48\textwidth}
    \begin{block}{Partie robot}
      \begin{itemize}
        \item r\'eception radio et d\'ecodage du paquet;
        \item commande diff\'erentielle des moteurs du Cutebot Pro;
        \item gestion des phares, du haut-parleur et des LEDs;
        \item mesure de distance par capteur ultrason.
      \end{itemize}
    \end{block}
  \end{columns}

  \vspace{0.4cm}
  \begin{beamercolorbox}[rounded=true,wd=\textwidth,sep=0.8em]{block body}
    \textbf{Organisation logicielle :} un seul projet Segger Embedded Studio,
    avec \texttt{main.c} qui s\'electionne \texttt{run\_tilt\_tx()} ou
    \texttt{run\_tilt\_rx()} selon la carte flash\'ee.
  \end{beamercolorbox}
\end{frame}

\begin{frame}{Architecture et fonctionnement}
  \begin{columns}[T,onlytextwidth]
    \column{0.6\textwidth}
    \begin{tikzpicture}[
      node distance=0.55cm and 0.55cm,
      box/.style={draw=Primary, rounded corners=3pt, fill=Soft, minimum width=3.0cm, minimum height=0.85cm, align=center},
      arr/.style={-{Latex[length=2.4mm]}, thick, draw=Primary}
    ]
      \node[box] (acc) {Acc\'el\'erom\`etre\\+ boutons};
      \node[box, below=of acc] (tx) {Mise \`a l'\'echelle\\et paquet radio};
      \node[box, below=of tx] (rx) {R\'eception\\sur le robot};
      \node[box, below left=0.7cm and 0.4cm of rx] (drive) {Calcul des vitesses\\gauche / droite};
      \node[box, below right=0.7cm and 0.4cm of rx] (safe) {V\'erification\\des obstacles};
      \node[box, below=1.1cm of rx] (act) {Moteurs + \\ phares + son};

      \draw[arr] (acc) -- (tx);
      \draw[arr] (tx) -- node[right, font=\scriptsize]{BLE LR} (rx);
      \draw[arr] (rx) -- (drive);
      \draw[arr] (rx) -- (safe);
      \draw[arr] (drive) -- (act);
      \draw[arr] (safe) -- (act);
    \end{tikzpicture}

    \column{0.36\textwidth}
    \begin{block}{Principe de commande}
      \begin{itemize}
        \item \texttt{ax} pilote la direction;
        \item \texttt{ay} pilote l'avance / recul;
        \item le robot calcule les deux vitesses roues:
      \end{itemize}
      \vspace{0.2cm}
      \centering
      \texttt{left = (ay - ax) / 2}\\
      \texttt{right = (ay + ax) / 2}
    \end{block}
  \end{columns}
\end{frame}

\begin{frame}{Fonctionnalit\'es principales}
  \begin{columns}[T,onlytextwidth]
    \column{0.48\textwidth}
    \begin{block}{Commande et interaction}
      \begin{itemize}
        \item pilotage par inclinaison;
        \item transmission radio des axes et des boutons;
        \item deux modes de stunt lanc\'es par A et B;
        \item retour visuel sur la matrice LED.
      \end{itemize}
    \end{block}

    \column{0.48\textwidth}
    \begin{block}{S\'ecurit\'e et retour utilisateur}
      \begin{itemize}
        \item \'evitement d'obstacle sous 30 cm;
        \item sortie d'\'evitement \`a 45 cm ou apr\`es timeout;
        \item bip p\'eriodique en marche arri\`ere.
      \end{itemize}
    \end{block}
  \end{columns}

  \vspace{0.4cm}
  \begin{alertblock}{Valeur ajout\'ee}
    Le projet ne se limite pas \`a d\'eplacer le robot : il int\`egre aussi
    la s\'ecurit\'e, les effets lumineux et sonores, et une machine \`a \'etats bien structur\'ee, assurant le bon fonctionnement de l'ensemble.
  \end{alertblock}
\end{frame}

\begin{frame}{Difficult\'es rencontr\'ees}
  \begin{columns}[T,onlytextwidth]
    \column{0.49\textwidth}
    \begin{block}{Difficult\'es techniques}
      \begin{itemize}
        \item comprendre et piloter plusieurs p\'eriph\'eriques en bas niveau;
        \item calibrer les axes pour obtenir un contr\^ole naturel;
        \item synchroniser radio, moteur, ultrason et retours utilisateur;
        \item g\'erer les erreurs de transmission et les cas de blocage.
      \end{itemize}
    \end{block}

    \column{0.49\textwidth}
    \begin{block}{R\'eponses apport\'ees}
      \begin{itemize}
        \item s\'eparation claire entre drivers et logique applicative;
        \item limitation des vitesses avec \texttt{clamp} et zone morte;
        \item machine d'\'etats pour les modes normal, \'evitement et stunt;
        \item temporisations et timeouts pour garder un comportement stable.
      \end{itemize}
    \end{block}
  \end{columns}

  \vspace{0.35cm}
\end{frame}

\begin{frame}{Bilan et perspectives}
  \begin{block}{Bilan}
    \begin{itemize}
      \item objectif principal atteint : le robot est contr\^ol\'e par inclinaison;
      \item la communication radio transmet mouvement et actions boutons;
      \item le syst\`eme r\'eagit aux obstacles;
      \item le code final est modulaire.
    \end{itemize}
  \end{block}



  \vspace{0.25cm}
  \centering
  \textbf{Conclusion :} le projet montre qu'un syst\`eme simple fonctionne tr\`es bien si l'on combine correctement contr\^ole, s\'ecurit\'e et code lisible.
\end{frame}

\end{document}
